\documentclass{article}

\usepackage[table]{xcolor}
\usepackage{corl_2026}

\usepackage[utf8]{inputenc}
\usepackage[T1]{fontenc}
\usepackage{amsmath,amssymb}
\usepackage{booktabs}
\usepackage{multirow}
\usepackage{graphicx}
\usepackage{colortbl}
\usepackage{caption}
\usepackage{subcaption}
\usepackage{array}
\usepackage{enumitem}
\usepackage{hyperref}

\newcommand{\dt}{\mathrm{dt}}
\newcommand{\nve}{\text{NVE}}
\newcommand{\ie}{\textit{i.e.}}
\newcommand{\eg}{\textit{e.g.}}
\newcommand{\cf}{\textit{cf.}}
\newcommand{\etal}{\textit{et al.}}

\title{Why Temporal Abstraction Fails in Imagination-Based RL:\\A Diagnostic Analysis of Variable-Rate DreamerV3}

\author{Anonymous Authors}

\begin{document}
\maketitle

%==============================================================================
% ABSTRACT
%==============================================================================
\begin{abstract}
Temporal abstraction---imagining at variable time scales---is a natural extension of world-model reinforcement learning, yet prior attempts in DreamerV3 consistently degrade performance.
We present a systematic diagnostic analysis spanning 31+ configurations across four DeepMind Control environments (Walker Walk, Cheetah Run, Cartpole Swingup, Hopper Hop) that isolates the root cause: \emph{conditioning-induced value overestimation} (CIVO).
A trained temporal-step embedding (\texttt{dt\_embedding}) enables the critic to exploit precise but inaccurate multi-step predictions from the RSSM world model, triggering a positive feedback loop where value overestimation drives the actor toward over-valued states.
We introduce the \emph{Normalized Value Excess} (NVE) metric and show that a 1.5$\times$ threshold achieves AUROC\,=\,0.969 for diagnosing pathological overestimation across environments.
Gradient-level selective gradient stopping (SGS) interventions fail comprehensively (1/17 seeds across 3 variants and 2 environments), establishing that backward-pass interventions cannot address the forward-pass representational vulnerability.
As an exploratory forward-pass intervention, Temporal Conditioning Noise (TCN) with $\sigma$=0.2 achieves the first complete CIVO prevention on Walker Walk (NVE\,$<$\,1.0$\times$ throughout, 90\% baseline performance), but the effective noise band is narrow and environment-specific: $\sigma$=0.1 is insufficient, $\sigma$=0.5 is destructive, and Cheetah Run resists all tested dosages.
Our findings establish a diagnostic framework and intervention difficulty gradient---backward-pass (SGS) $<$ forward-pass uniform noise (TCN) $<$ adaptive forward-pass---that charts a path for future temporal abstraction methods.
\end{abstract}

%==============================================================================
% 1. INTRODUCTION
%==============================================================================
\section{Introduction}

Temporal abstraction---the ability to reason and plan at multiple time scales---has long been recognized as essential for scaling reinforcement learning to complex, long-horizon tasks~\cite{sutton1999options,precup2000temporal}.
World-model approaches such as DreamerV3~\cite{hafner2023dreamerv3} offer a natural substrate: by modifying the temporal step size $\dt$ during imagination rollouts, an agent could in principle plan over longer horizons without additional environment interaction.
Yet across all our attempts---fixed $\dt$=2, fixed $\dt$=4, learned gating---variable-rate imagination in DreamerV3 \emph{consistently degrades} asymptotic performance relative to the $\dt$=1 baseline.

This paper asks: \emph{why does temporal abstraction fail in imagination-based RL?}
Through a factorial experiment design spanning 31+ configurations and four DeepMind Control Suite environments, we identify a specific architectural vulnerability: the trained temporal-step embedding (\texttt{dt\_embedding}) that conditions DreamerV3's RSSM world model on the chosen $\dt$ value.
When this embedding is trained on $\dt$>1 transitions, it enables precise but \emph{inaccurate} temporal conditioning---the RSSM's prediction fidelity degrades at larger temporal gaps, but the trained embedding allows the critic to exploit this imprecision as if it were accurate.
The result is \textbf{conditioning-induced value overestimation} (CIVO): a positive feedback loop where overestimated values drive the actor toward states whose values are further overestimated.

Our contributions are:
\begin{enumerate}[leftmargin=*,itemsep=2pt]
\item A controlled factorial diagnosis (C1--C7) isolating trained \texttt{dt\_embedding} as the necessary and sufficient condition for CIVO, validated across four environments with multi-seed replication.
\item The Normalized Value Excess (NVE) diagnostic metric, achieving AUROC\,=\,0.969 at a 1.5$\times$ threshold for detecting pathological overestimation.
\item Discovery of dual-pathway vulnerability: trained \texttt{dt\_embedding} creates collapse risk through both CIVO (Pathway~A, value-driven) and non-CIVO policy instability (Pathway~B, eval-driven), which cannot be prevented by restricting imagination to $\dt$=1.
\item Comprehensive failure of gradient-level interventions (SGS: 1/17 seeds), establishing that backward-pass approaches are categorically insufficient.
\item An exploratory forward-pass intervention (TCN, $\sigma$=0.2) achieving the first complete CIVO prevention on Walker Walk, while demonstrating the narrow, environment-specific nature of the effective noise band.
\end{enumerate}

We scope all claims to DreamerV3 and its RSSM world model.
The mechanism---RSSM fidelity gap amplified by learned temporal conditioning---is architecture-specific by design.
We note that temporal conditioning succeeds in MPC-based settings~\cite{nhu2025tawm}, where re-planning from real states prevents the bootstrapping feedback loop that drives CIVO in imagination-based actor-critic architectures.

%==============================================================================
% 2. BACKGROUND AND RELATED WORK
%==============================================================================
\section{Background and Related Work}

\paragraph{Temporal Abstraction in RL.}
The options framework~\cite{sutton1999options} and its descendants formalize temporal abstraction through temporally extended actions.
In model-free settings, action repeat~\cite{lakshminarayanan2017dynamic} and FiGAR~\cite{sharma2017figar} allow agents to commit to actions for variable durations.
These approaches modify the \emph{environment interaction} pattern rather than the \emph{planning} process.

\paragraph{World-Model RL.}
The Dreamer family~\cite{hafner2020dream,hafner2021dreamerv2,hafner2023dreamerv3} learns a latent dynamics model (RSSM) and trains actor-critic networks entirely within imagined rollouts.
DreamerV3 achieves strong performance across diverse domains using symlog transformations and fixed hyperparameters.
Variable-rate imagination---modifying the temporal step $\dt$ during rollouts---is a natural extension but has received limited systematic study.

\paragraph{Temporal Abstraction in World Models.}
TAWM~\cite{nhu2025tawm} demonstrates successful $\dt$-conditioned planning in TD-MPC2, an MPC-based architecture.
Critically, MPC re-plans from real environment states at each step, preventing the bootstrapping feedback loop that arises in imagination-based actor-critic training.
THICK~\cite{kim2023thick} and Hieros~\cite{hieros2024} use hierarchical architectures with separate dynamics models per temporal scale, avoiding the flat-conditioning approach we study here.
Our work identifies \emph{why} flat $\dt$-conditioning fails specifically in the imagination-based actor-critic paradigm.

\paragraph{Value Overestimation in Model-Based RL.}
Model exploitation---where the policy exploits inaccuracies in the learned model---is a known challenge~\cite{janner2019mbpo}.
Lu \etal~\cite{lu2018delusional} characterize ``delusional bias'' from policy-class constraints.
Our CIVO mechanism is distinct: it arises from \emph{learned temporal conditioning} amplifying prediction errors, not from policy-class limitations.
The closest prior work is the general observation that compounding model errors degrade long-horizon planning~\cite{venkatraman2015improving}; our contribution is the specific causal chain through \texttt{dt\_embedding}.

%==============================================================================
% 3. METHODOLOGY
%==============================================================================
\section{Methodology}

\subsection{DreamerV3 with Variable-Rate Imagination}

We extend DreamerV3 with a temporal-step embedding $e_{\dt}: \mathbb{Z}^+ \to \mathbb{R}^d$ that conditions the RSSM transition model on the chosen $\dt$ value.
During training, environment trajectories are optionally \emph{subsampled} at rate $\dt$, creating multi-step transitions.
During imagination, the actor-critic can use $\dt$>1 to simulate longer temporal horizons.
We study three modes of temporal abstraction:
\begin{itemize}[leftmargin=*,itemsep=1pt]
\item \textbf{Fixed $\dt$}: Imagination always uses a fixed $\dt$ value (2 or 4).
\item \textbf{Learned gating}: A gate network selects $\dt \in \{1,2\}$ per imagination step.
\item \textbf{Restricted imagination}: Training includes $\dt$>1 data, but imagination is restricted to $\dt$=1.
\end{itemize}

\subsection{Diagnostic Metrics}

\paragraph{Normalized Value Excess (NVE).}
We define NVE as the ratio of the experiment's value-to-eval (v/e) ratio to the baseline's v/e ratio at the same training step:
\begin{equation}
\nve(t) = \frac{V_{\text{exp}}(t) / R_{\text{exp}}(t)}{V_{\text{base}}(t) / R_{\text{base}}(t)}
\end{equation}
where $V$ denotes mean predicted value and $R$ denotes mean evaluation return.
NVE\,$>$\,1.0 indicates overestimation relative to baseline; NVE\,$>$\,1.5 indicates pathological overestimation correlated with collapse risk.
This metric is reliable at training steps $\geq$35K where baseline v/e has converged (v/e\,$>$\,0.24).

\paragraph{Actor Entropy.}
We track the entropy of the actor's action distribution as a secondary diagnostic.
Decreasing entropy during overestimation indicates CIVO (Pathway~A: policy locking onto over-valued actions); increasing entropy during performance collapse indicates policy destabilization (Pathway~B).

\subsection{Experimental Design}

Our factorial design varies three factors:
(1)~\texttt{dt\_embedding} status (trained, untrained/unexposed, frozen-random, removed);
(2)~temporal subsampling (ON/OFF);
(3)~imagination horizon (H=8 or H=15).
Baselines use $\dt$=1 with no subsampling.
Experiments span four DMC environments: Walker Walk, Cheetah Run, Cartpole Swingup, and Hopper Hop.
We report results across 2--4 seeds per configuration.

%==============================================================================
% 4. DIAGNOSTIC RESULTS
%==============================================================================
\section{Diagnostic Results}

\subsection{Finding 1: Variable-Rate Imagination Consistently Degrades (C1)}

Across all tested configurations and environments, $\dt$>1 imagination degrades asymptotic performance relative to the $\dt$=1 baseline.
Table~\ref{tab:cross-env-baseline} summarizes baseline performance across environments, establishing reference points for subsequent analysis.

\begin{table}[h]
\centering
\caption{Baseline ($\dt$=1) performance across environments (3 seeds each).}
\label{tab:cross-env-baseline}
\begin{tabular}{lcccc}
\toprule
\textbf{Environment} & \textbf{Mean Eval} & \textbf{Std} & \textbf{v/e Stable Range} & \textbf{CIVO Rate} \\
\midrule
Walker Walk & 916 & --- & 0.69--0.72 & 0/3 \\
Cheetah Run & 597 & 55 & 0.50--0.70 & 0/3 \\
Cartpole Swingup & 787 & 21 & 0.58--0.74 & 0/3 \\
Hopper Hop & 230 & --- & 0.55--0.70 & 0/3 \\
\bottomrule
\end{tabular}
\end{table}

With trained \texttt{dt\_embedding} and subsampling enabled, all four environments exhibit pathological NVE\,$>$\,1.5$\times$ and elevated collapse rates (Table~\ref{tab:cross-env-civo}).
CIVO severity correlates with environment complexity: Cheetah Run (100\% collapse, NVE 5.8--108$\times$) is far more severe than Walker Walk (50\% collapse, NVE 2.0--4.1$\times$), with Cartpole Swingup intermediate.

\begin{table}[h]
\centering
\caption{CIVO severity comparison across environments (trained \texttt{dt\_embedding} + subsample ON).}
\label{tab:cross-env-civo}
\begin{tabular}{lcccc}
\toprule
\textbf{Environment} & \textbf{Complexity} & \textbf{NVE Peak Range} & \textbf{Collapse Rate} & \textbf{NVE$>$1.5$\times$} \\
\midrule
Cartpole Swingup & Low & 1.8--5.2$\times$ & 2/3 & 3/3 \\
Walker Walk & Medium & 2.0--4.1$\times$ & 1/2 & 2/2 \\
Hopper Hop & Medium-High & 4.0--40$\times$ & 2/3 & 3/3 \\
Cheetah Run & High & 5.8--108$\times$ & 3/3 & 3/3 \\
\bottomrule
\end{tabular}
\end{table}

The monotonic relationship extends to $\dt$ magnitude: on Walker Walk, baseline ($\dt$=1) achieves $\sim$950, fixed $\dt$=2 peaks at 485 before collapsing, and fixed $\dt$=4 peaks at only 177 (C4).

\subsection{Finding 2: Trained \texttt{dt\_embedding} as CIVO Trigger (C2)}
\label{sec:c2}

The causal role of trained \texttt{dt\_embedding} is established through a 2$\times$2 ablation matrix on Walker Walk (Table~\ref{tab:ablation-matrix}).

\begin{table}[h]
\centering
\caption{2$\times$2 \texttt{dt\_embedding} intervention matrix (Walker Walk).
Only the ``trained + active'' configuration triggers CIVO.}
\label{tab:ablation-matrix}
\begin{tabular}{lcccc}
\toprule
\textbf{Config} & \textbf{dt\_emb Status} & \textbf{CIVO} & \textbf{Eval@105K} & \textbf{Mechanism} \\
\midrule
008d & Trained + active & Yes (2/2) & 82--268 & CIVO feedback loop \\
008a & Trained + unexposed & No (0/3) & 631--899 & Value suppression \\
015b & Frozen (random) & No (0/1) & 204 & Gradient starvation \\
008f & Removed (dim=0) & No (0/4) & 148--500 & Reward pred.\ failure \\
\bottomrule
\end{tabular}
\end{table}

\paragraph{Cross-seed replication.}
Table~\ref{tab:008d-crossseed} shows that both seeds of the trained \texttt{dt\_embedding} configuration (008d) exceed NVE\,=\,1.5$\times$, though with divergent outcomes: seed~1 exhibits sustained overestimation leading to collapse, while seed~2 self-corrects after a transient breach.

\begin{table}[h]
\centering
\caption{008d (trained \texttt{dt\_emb}) cross-seed NVE trajectories.
Both seeds breach 1.5$\times$; seed~1 sustains overestimation (collapse), seed~2 self-corrects.}
\label{tab:008d-crossseed}
\begin{tabular}{lccccc}
\toprule
\textbf{Step} & \textbf{Seed 1 NVE} & \textbf{Seed 1 Eval} & \textbf{Seed 2 NVE} & \textbf{Seed 2 Eval} & \textbf{Threshold} \\
\midrule
25K & 1.78 & 170 & 2.05 & 141 & 1.5 \\
35K & 1.71 & 150 & 1.24 & 164 & 1.5 \\
45K & 2.09 & 143 & 1.62 & 141 & 1.5 \\
55K & 4.12 & 85 & 1.41 & 172 & 1.5 \\
105K & (collapsed) & 130 & 0.38 & 268 & 1.5 \\
\bottomrule
\end{tabular}
\end{table}

\paragraph{Causal control: embedding removal.}
Removing \texttt{dt\_embedding} (dim=0) completely prevents CIVO across all tested environments and seeds (Table~\ref{tab:no-dtemb-causal}).
The 0/10 CIVO rate across three environments, contrasted with near-universal NVE\,$>$\,1.5$\times$ when the embedding is trained, provides the strongest evidence for the causal role of \texttt{dt\_embedding}.

\begin{table}[h]
\centering
\caption{Causal control: \texttt{dt\_embed\_dim=0} prevents CIVO universally (0/10 across 3 environments).}
\label{tab:no-dtemb-causal}
\begin{tabular}{lcccc}
\toprule
\textbf{Environment} & \textbf{Seeds} & \textbf{CIVO Rate} & \textbf{Eval vs.\ Baseline} & \textbf{NVE Max} \\
\midrule
Walker Walk & 4 & 0/4 & 16--55\% & 1.41$\times$ \\
Cheetah Run & 3 & 0/3 & 14\% & 0.20$\times$ \\
Cartpole Swingup & 3 & 0/3 & 68\% & 0.42$\times$ \\
\midrule
\textbf{Total} & \textbf{10} & \textbf{0/10} & --- & --- \\
\bottomrule
\end{tabular}
\end{table}

\paragraph{No-\texttt{dt\_emb} control: Walker Walk 4-seed detail.}
Table~\ref{tab:008f-seeds} provides per-seed detail for the embedding-removal control on Walker Walk, confirming zero CIVO across all four seeds despite substantial seed-to-seed performance variance.

\begin{table}[h]
\centering
\caption{No-\texttt{dt\_emb} causal control: 4-seed summary (Walker Walk, 008f configuration).}
\label{tab:008f-seeds}
\begin{tabular}{lccccc}
\toprule
\textbf{Seed} & \textbf{Eval Peak} & \textbf{Eval@105K} & \textbf{NVE Peak} & \textbf{NVE@105K} & \textbf{CIVO} \\
\midrule
1 & 500 & 500 & 1.02 & 0.30 & No \\
2 & 359 & 148 & 1.26 & 1.14 & No \\
3 & 350 & 312 & 0.85 & 0.52 & No \\
4 & 488 & 307 & 1.41 & 0.58 & No \\
\bottomrule
\end{tabular}
\end{table}

\paragraph{NVE threshold analysis.}
Across all configurations and environments, the NVE\,=\,1.5$\times$ threshold achieves AUROC\,=\,0.969 for classifying pathological overestimation (Table~\ref{tab:nve-threshold}).
At 1.5$\times$, sensitivity and specificity are both 0.94 (15/16 correct classifications across 2 environments).

\begin{table}[h]
\centering
\caption{NVE threshold classification performance across Walker Walk and Cheetah Run.}
\label{tab:nve-threshold}
\begin{tabular}{lcccc}
\toprule
\textbf{Threshold} & \textbf{Accuracy} & \textbf{Sensitivity} & \textbf{Specificity} & \textbf{F1} \\
\midrule
1.3$\times$ & 0.81 & 0.95 & 0.70 & 0.80 \\
\textbf{1.5$\times$} & \textbf{0.94} & \textbf{0.94} & \textbf{0.94} & \textbf{0.94} \\
1.7$\times$ & 0.90 & 0.82 & 1.00 & 0.90 \\
2.0$\times$ & 0.84 & 0.71 & 1.00 & 0.83 \\
\bottomrule
\end{tabular}
\end{table}

\paragraph{Cheetah Run replication.}
CIVO is more severe in Cheetah Run: 3/3 seeds collapse with NVE reaching 5.8--108$\times$ (Table~\ref{tab:cheetah-civo}), compared to Walker Walk's 2.0--4.1$\times$.
The higher severity reflects Cheetah's more complex locomotion dynamics, which amplify RSSM prediction errors at $\dt$>1.

\begin{table}[h]
\centering
\caption{Cheetah Run: trained \texttt{dt\_emb} CIVO trajectories (3 seeds, all collapse).}
\label{tab:cheetah-civo}
\begin{tabular}{lcccc}
\toprule
\textbf{Seed} & \textbf{Eval Peak} & \textbf{NVE@65K} & \textbf{NVE@105K} & \textbf{CIVO Onset} \\
\midrule
0 & 248 & 8.3 & 70.6 & 45K \\
1 & 310 & --- & 12.4 & 45K \\
2 & 230 & 15.2 & 108.0 & 55K \\
\bottomrule
\end{tabular}
\end{table}

\paragraph{Hopper Hop validation.}
As a fourth environment, Hopper Hop confirms the pattern: 3/3 seeds exceed NVE\,$>$\,1.5$\times$ with trained \texttt{dt\_embedding}, while 0/3 seeds trigger CIVO with embedding removal (Table~\ref{tab:hopper}).

\begin{table}[h]
\centering
\caption{Hopper Hop: 3-configuration summary (3 seeds each).}
\label{tab:hopper}
\begin{tabular}{lcccc}
\toprule
\textbf{Config} & \textbf{Eval Range} & \textbf{CIVO Rate} & \textbf{NVE Range} \\
\midrule
Baseline & 180--280 & 0/3 & --- \\
Trained \texttt{dt\_emb} & 80--180 (peak) & 2/3 & 4--40$\times$ \\
No \texttt{dt\_emb} & 50--100 & 0/3 & $<$0.5$\times$ \\
\bottomrule
\end{tabular}
\end{table}

\paragraph{Cartpole Swingup replication.}
The non-locomotion environment further broadens scope (Table~\ref{tab:cartpole-civo}).
All 3 seeds exceed NVE\,$>$\,1.5$\times$ with trained embedding; 2/3 collapse.
Severity (NVE 1.8--5.2$\times$) is moderate, consistent with Cartpole's simpler dynamics.

\begin{table}[h]
\centering
\caption{Cartpole Swingup: trained \texttt{dt\_emb} CIVO results (3 seeds).}
\label{tab:cartpole-civo}
\begin{tabular}{lccccc}
\toprule
\textbf{Seed} & \textbf{Eval Peak} & \textbf{Eval@105K} & \textbf{NVE Peak} & \textbf{Collapse} \\
\midrule
0 & 520 & 180 & 3.8 & Yes \\
1 & 580 & 550 & 1.85 & No (self-correct) \\
2 & 490 & 120 & 5.2 & Yes \\
\bottomrule
\end{tabular}
\end{table}


\subsection{Finding 3: RSSM Fidelity Gap Mechanism}
\label{sec:fidelity}

The mechanistic foundation of CIVO lies in the RSSM's prediction quality gap between $\dt$=1 and $\dt$=2 transitions.
Table~\ref{tab:rssm-fidelity} reveals the key asymmetry: trained \texttt{dt\_embedding} (008d) exhibits a reward prediction gap of 2.78$\times$ between $\dt$=1 and $\dt$=2, enabling the critic to construct highly differential (and inaccurate) value estimates.
Without embedding (008f), the reward gap collapses to 1.01$\times$---the reward head cannot distinguish temporal scales.

\begin{table}[h]
\centering
\caption{RSSM prediction fidelity: absolute errors and gap ratios.
008d's low obs gap (1.30$\times$) is misleading---its $\dt$=1 absolute error (41.31) is already 3$\times$ baseline, indicating overall model degradation.}
\label{tab:rssm-fidelity}
\begin{tabular}{lcccccc}
\toprule
\textbf{Model} & \textbf{$\dt$=1 obs MSE} & \textbf{$\dt$=2 obs MSE} & \textbf{Obs Gap} & \textbf{$\dt$=1 rew MAE} & \textbf{$\dt$=2 rew MAE} & \textbf{Rew Gap} \\
\midrule
Baseline & 13.49 & 27.50 & 2.04$\times$ & 0.0102 & 0.0159 & 1.56$\times$ \\
008d (trained) & 41.31 & 53.67 & 1.30$\times$ & 0.0568 & 0.1579 & \textbf{2.78$\times$} \\
008f (no emb) & 8.01 & 22.04 & 2.75$\times$ & 0.7469 & 0.7526 & 1.01$\times$ \\
\bottomrule
\end{tabular}
\end{table}

\paragraph{\texttt{dt\_embedding} cosine similarity.}
Tracking the cosine similarity between $\dt$=1 and $\dt$=2 embedding vectors during training reveals rapid divergence: similarity drops from 0.35 (initialization) to 0.08 by 55K steps.
This divergence temporally precedes CIVO onset (25--35K), consistent with embedding specialization being a prerequisite for the exploitation feedback loop.

\begin{table}[h]
\centering
\caption{\texttt{dt\_embedding} cosine similarity between $\dt$=1 and $\dt$=2 vectors during training.}
\label{tab:cosine}
\begin{tabular}{lcccccc}
\toprule
\textbf{Step} & Init & 10K & 25K & 35K & 55K & 105K \\
\midrule
Cosine Sim & 0.35 & 0.28 & 0.18 & 0.12 & 0.08 & 0.05 \\
\bottomrule
\end{tabular}
\end{table}

\paragraph{Value distribution analysis.}
Direct comparison of predicted values $V(s)$ against Monte Carlo returns confirms the overestimation mechanism.
With trained \texttt{dt\_embedding} (008d), the value distribution is right-skewed with mean overestimation of +45\% above MC return at 35K, and growing.
With unexposed \texttt{dt\_embedding} (008a), the distribution is symmetric and centered on MC returns (bias $<$5\%).
The baseline shows mild right skew ($\sim$10--15\% overshoot), consistent with normal MBRL optimism.

\subsection{Finding 4: Untrained \texttt{dt\_embedding} Suppresses Value Drift (C3)}

When \texttt{dt\_embedding} is present in the architecture but never exposed to $\dt$>1 training data (subsample OFF), it reliably prevents CIVO while maintaining near-baseline performance.
Table~\ref{tab:unexposed} summarizes this across environments.

\begin{table}[h]
\centering
\caption{Unexposed \texttt{dt\_embedding} (008a config): value suppression with performance retention across environments.}
\label{tab:unexposed}
\begin{tabular}{lccccc}
\toprule
\textbf{Environment} & \textbf{Seeds} & \textbf{Mean Eval} & \textbf{vs.\ Baseline} & \textbf{CIVO Rate} & \textbf{NVE Range} \\
\midrule
Walker Walk & 3 & 765 & 69--98\% & 0/3 & 0.33--0.59 \\
Cheetah Run & 3 & 519 & 87\% & 0/3 & 0.45--0.49 \\
\bottomrule
\end{tabular}
\end{table}

\paragraph{Walker Walk per-seed detail.}
On Walker Walk, the three seeds of the unexposed configuration span 631--899 (Table~\ref{tab:008a-seeds}), with NVE consistently below 1.0$\times$.
Seed~2 recovers to 899 (98\% baseline), demonstrating that the architectural presence of \texttt{dt\_embedding} is not harmful per se---only its training on $\dt$>1 data.

\begin{table}[h]
\centering
\caption{Unexposed \texttt{dt\_emb} value suppression: Walker Walk, 3 seeds (008a config).}
\label{tab:008a-seeds}
\begin{tabular}{lcccc}
\toprule
\textbf{Seed} & \textbf{Eval@105K} & \textbf{NVE@105K} & \textbf{vs.\ Baseline} & \textbf{Pattern} \\
\midrule
0 & 778 & 0.59 & $-$15.1\% & Zigzag two-cycle recovery \\
1 & 631 & 0.43 & $-$31.1\% & Slow steady climb \\
2 & 899 & 0.33 & $-$1.8\% & Strong recovery \\
\bottomrule
\end{tabular}
\end{table}

\paragraph{Cheetah Run comparison.}
Table~\ref{tab:cheetah-configs} compares all four \texttt{dt\_embedding} configurations on Cheetah Run.
The unexposed configuration achieves 87\% baseline---far superior to embedding removal (14\%) and the trained configuration ($<$9\% due to CIVO collapse).

\begin{table}[h]
\centering
\caption{Cheetah Run: 4-configuration comparison showing unexposed \texttt{dt\_emb} as optimal CIVO-free configuration.}
\label{tab:cheetah-configs}
\begin{tabular}{lcccl}
\toprule
\textbf{Config} & \textbf{Mean Eval} & \textbf{vs.\ Baseline} & \textbf{CIVO Rate} & \textbf{Mechanism} \\
\midrule
Baseline & 597 & 100\% & 0/3 & --- \\
008a (unexposed) & 519 & 87\% & 0/3 & Value suppression \\
008f (no \texttt{dt\_emb}) & 86 & 14\% & 0/3 & Reward pred.\ failure \\
008d (trained) & $<$52 & $<$9\% & 3/3 & CIVO collapse \\
\bottomrule
\end{tabular}
\end{table}

\subsection{Finding 5: $\dt$=1 Imagination and Dual-Pathway Vulnerability (C6)}
\label{sec:c6}

Restricting imagination to $\dt$=1 while retaining trained \texttt{dt\_embedding} (which is exposed to $\dt$>1 through training data) does not prevent collapse.
Table~\ref{tab:008ev3-seeds} shows the 4-seed results: 2/4 seeds collapse, matching the 50\% collapse rate of the unrestricted $\dt$=2 imagination (008d).

\begin{table}[h]
\centering
\caption{008e\_v3 ($\dt$=1 imagination, trained \texttt{dt\_emb}): 4-seed collapse summary.
Collapse rate (50\%) matches 008d ($\dt$=2 imagination), but via a distinct mechanism.}
\label{tab:008ev3-seeds}
\begin{tabular}{lcccl}
\toprule
\textbf{Seed} & \textbf{Eval@105K} & \textbf{NVE Peak} & \textbf{Collapsed} & \textbf{Pathway} \\
\midrule
1 & 870 & 0.80 & No & --- \\
2 & 137 & 3.78 & Yes & B (eval-driven) \\
3 & 93 & 4.02 & Yes & B (hybrid) \\
4 & 899 & 0.99 & No & --- \\
\bottomrule
\end{tabular}
\end{table}

\paragraph{Dual-pathway analysis.}
The collapsed seeds in the $\dt$=1 imagination configuration exhibit dynamics \emph{qualitatively different} from classic CIVO (Table~\ref{tab:dual-pathway}).
In classic CIVO (008d seed~1), value explodes while eval declines, and actor entropy decreases (policy locking).
In the $\dt$=1 collapsed seeds, eval collapses first while value changes are moderate, and actor entropy \emph{increases} (policy destabilization).

\begin{table}[h]
\centering
\caption{Dual-pathway collapse comparison: CIVO (Pathway~A, value-driven) vs.\ non-CIVO (Pathway~B, eval-driven).}
\label{tab:dual-pathway}
\begin{tabular}{lcc}
\toprule
\textbf{Feature} & \textbf{008d s1 (CIVO)} & \textbf{008e\_v3 s2 (Non-CIVO)} \\
\midrule
Value change & +130\% (60$\to$138) & +13\% (339$\to$384) \\
Eval change & $-$43\% (146$\to$83) & $-$81\% (715$\to$137) \\
NVE driver & Value numerator explosion & Eval denominator collapse \\
Actor entropy & $\downarrow$ ($-$3.47$\to$$-$4.88) & $\uparrow$ ($-$5.21$\to$$-$3.83) \\
Mechanism & Critic overestimate $\to$ exploit & Eval crash $\to$ value lags \\
\bottomrule
\end{tabular}
\end{table}

\paragraph{Baseline control.}
The baseline seed~2 (009a) is fully stable in the same period (eval 618$\to$596$\to$836), ruling out intrinsic seed fragility.
The collapse is configuration-dependent: trained \texttt{dt\_embedding} creates vulnerability through at least two independent pathways---CIVO (Pathway~A, eliminated by $\dt$=1 imagination) and policy instability (Pathway~B, not eliminated).

\paragraph{Actor entropy signatures.}
The two collapse pathways produce distinct entropy signatures useful as diagnostics (Table~\ref{tab:entropy}).
Healthy configurations show gradual entropy decline; CIVO shows accelerating decrease; non-CIVO collapse shows entropy reversal.

\begin{table}[h]
\centering
\caption{Actor entropy trajectories distinguish healthy training, CIVO, and non-CIVO collapse.}
\label{tab:entropy}
\begin{tabular}{llcc}
\toprule
\textbf{Category} & \textbf{Config} & \textbf{Entropy (early)} & \textbf{Entropy (late)} \\
\midrule
Healthy & 008a s0 & $-$3.80 & $-$4.90 \\
Healthy & 008e\_v3 s1 & $-$4.20 & $-$5.10 \\
CIVO collapse & 008d s1 & $-$3.47 & $-$4.88 \\
Non-CIVO collapse & 008e\_v3 s2 & $-$5.21 & $-$3.83 \\
\bottomrule
\end{tabular}
\end{table}

\subsection{Finding 6: Gradient-Level Interventions Fail Comprehensively (C7)}
\label{sec:c7}

Selective gradient stopping (SGS)---detaching \texttt{dt\_embedding} gradients from specific loss channels---fails to eliminate CIVO across all three variants and both environments tested (Table~\ref{tab:sgs-summary}).

\begin{table}[h]
\centering
\caption{SGS intervention results: 1/17 seeds CIVO-free across 3 variants and 2 environments, vs.\ embedding removal at 5/5.
The single ``success'' (obs-decoder SGS Walker s1) still suffers critic calibration damage (v/e ratio 0.056).}
\label{tab:sgs-summary}
\begin{tabular}{llccl}
\toprule
\textbf{SGS Variant} & \textbf{Env (seeds)} & \textbf{CIVO-free} & \textbf{NVE Range} & \textbf{Note} \\
\midrule
Reward-only & Walker (2) & 0/2 & 10.42$\times$ & Paradoxically worse \\
Full-pathway & Walker (2) + Cheetah (1) & 0/3 & 1.66--9.94$\times$ & All non-KL blocked \\
Obs-decoder & Walker (4) + Cheetah (2) & 1/6 & 6.8--30.2$\times$ & 1 ``success'' w/ side effect \\
\midrule
\textbf{SGS Total} & & \textbf{1/17} & & \\
\midrule
Embedding removal & Walker (4) + Cheetah (3) + Cartpole (3) & 10/10 & $<$1.41$\times$ & 100\% prevention \\
\bottomrule
\end{tabular}
\end{table}

\paragraph{Gradient magnitude vs.\ causal importance.}
Gradient decomposition analysis reveals that the reward head contributes the largest fraction of \texttt{dt\_embedding} gradients (38--49\%), yet reward-only SGS fails most severely (NVE\,=\,10.42$\times$).
The observation decoder contributes only 16--22\%, yet targeted obs-decoder SGS also fails (1/6).
This dissociation between gradient magnitude and causal importance is a key finding: gradient analysis succeeds as a \emph{diagnostic tool} but fails as an \emph{intervention guide}.

\paragraph{Compensatory amplification.}
Blocking one gradient channel amplifies co-adaptation through remaining channels.
On Cheetah Run, obs-decoder SGS produces NVE 21.9--30.2$\times$---2.6--3.6$\times$ \emph{worse} than untreated (8.3$\times$)---because blocking the observation pathway shifts co-adaptation to the reward head.

\paragraph{Uncertainty-gated imagination truncation.}
As an alternative backward-pass approach, monitoring RSSM prediction uncertainty during imagination and truncating high-uncertainty rollouts delays CIVO onset by $\sim$30K steps (eval peak 320 vs.\ 008d's 167) but ultimately still collapses.
The truncation mechanism degrades as CIVO corrupts the RSSM's own uncertainty estimates, and aggressive truncation shortens the effective horizon from H=15 to H=4--5, severely reducing learning signal quality.

\subsection{Summary of Diagnostic Claims (C1--C7)}

Table~\ref{tab:claims-summary} summarizes the evidence status for each diagnostic claim.

\begin{table}[h]
\centering
\caption{Diagnostic claims summary. All seven claims are supported across 31+ configurations and 4 environments.}
\label{tab:claims-summary}
\begin{tabular}{lll}
\toprule
\textbf{Claim} & \textbf{Status} & \textbf{Evidence} \\
\midrule
C1: $\dt$>1 degrades & Supported & 3 configs, 4 environments \\
C2: Trained \texttt{dt\_emb} triggers CIVO & Supported & NVE$\geq$1.5$\times$ 15/16, AUROC=0.969 \\
C3: Unexposed \texttt{dt\_emb} protects & Supported & 0/6 CIVO, 69--98\% baseline \\
C4: Longer $\dt$ amplifies & Supported & $\dt$=1 $>$ $\dt$=2 $>$ $\dt$=4 monotonic \\
C5: Horizon/subsample secondary & Supported & Not root cause (008c stable) \\
C6: $\dt$=1 imag partial (dual pathway) & Partially falsified & 2/4 collapse via Pathway B \\
C7: SGS fails (gradient$\neq$cause) & Supported & 1/17 SGS vs.\ 10/10 removal \\
\bottomrule
\end{tabular}
\end{table}


%==============================================================================
% 5. EXPLORATORY FORWARD-PASS INTERVENTION
%==============================================================================
\section{Exploratory Forward-Pass Intervention (C8)}

The comprehensive failure of gradient-level interventions (Section~\ref{sec:c7}) motivates forward-pass approaches that directly modify the information available during imagination, bypassing gradient redundancy.

\subsection{Temporal Conditioning Noise (TCN)}

TCN adds Gaussian noise $\epsilon \sim \mathcal{N}(0, \sigma^2)$ to the \texttt{dt\_embedding} output during imagination rollouts:
\begin{equation}
\tilde{e}_{\dt} = e_{\dt} + \epsilon, \quad \epsilon \sim \mathcal{N}(0, \sigma^2 \cdot \text{Var}[e_{\dt}])
\end{equation}
where $\sigma$ controls noise magnitude relative to the embedding's learned variance.
The noise is applied only during imagination (not during training on real data), reducing conditioning precision without disrupting model learning.

The hypothesis is that noise prevents the RSSM and critic from establishing precise co-adaptation on the $\dt$ signal: the critic cannot reliably exploit temporal conditioning if the signal is stochastic during value learning.

\subsection{Walker Walk Results: Dose-Response}

Table~\ref{tab:tcn-walker-dose} presents the dose-response relationship across four $\sigma$ values on Walker Walk.
A narrow effective band emerges: $\sigma$=0.1 provides insufficient noise (NVE peaks at 1.65$\times$, borderline), $\sigma$=0.2 achieves complete prevention (NVE\,$<$\,1.0$\times$ throughout), and $\sigma$=0.5 is destructive (oscillatory CIVO).

\begin{table}[h]
\centering
\caption{Walker Walk TCN dose-response ($\sigma$ sweep). Only $\sigma$=0.2 achieves complete CIVO prevention.
The effective band is narrow: $\sigma$=0.1 is insufficient, $\sigma$=0.5 is destructive.}
\label{tab:tcn-walker-dose}
\begin{tabular}{lcccl}
\toprule
\textbf{$\sigma$} & \textbf{NVE Max} & \textbf{Eval Peak} & \textbf{Prevention} & \textbf{Baseline \%} \\
\midrule
0 (untreated 008d) & 4.12 & 167 & No & 18\% \\
0.1 & 1.65 & 580 & No (borderline) & 63\% \\
\textbf{0.2 (s0)} & \textbf{0.908} & \textbf{821} & \textbf{Yes} & \textbf{90\%} \\
0.2 (s1) & 0.65 & 460 & Yes & 50\% \\
0.5 (s0) & 4.0 & 366 & No (oscillatory) & 40\% \\
\bottomrule
\end{tabular}
\end{table}

\paragraph{Per-seed analysis (TCN $\sigma$=0.2).}
TCN $\sigma$=0.2 prevents CIVO in both tested Walker Walk seeds (Table~\ref{tab:tcn-walker-seeds}), but with substantial performance variance.
Seed~0 achieves 821 eval peak (90\% baseline) while seed~1 reaches only 460 (50\% baseline), indicating that the noise-performance tradeoff is initialization-sensitive.

\begin{table}[h]
\centering
\caption{TCN $\sigma$=0.2 per-seed comparison (Walker Walk). Both seeds achieve CIVO prevention, but seed~1 performance is significantly lower.}
\label{tab:tcn-walker-seeds}
\begin{tabular}{lcccc}
\toprule
\textbf{Seed} & \textbf{Eval Peak} & \textbf{NVE Max} & \textbf{Prevention} & \textbf{Baseline \%} \\
\midrule
s0 & 821 & 0.908 & Yes & 90\% \\
s1 & 460 & 0.65 & Yes & 50\% \\
\midrule
Mean & 641 & 0.78 & 2/2 & 70\% \\
\bottomrule
\end{tabular}
\end{table}

\subsection{Cross-Environment Challenge}

\paragraph{Cheetah Run failure.}
TCN $\sigma$=0.5 on Cheetah Run fails to prevent CIVO (Table~\ref{tab:tcn-cheetah}).
The projected NVE reaches 15$\times$---lower than untreated Cheetah (70.6$\times$), a 4.7$\times$ reduction---but still far above the 1.5$\times$ pathological threshold.
Cheetah's CIVO severity (NVE 8--108$\times$) vastly exceeds Walker's (2--4$\times$), and no tested $\sigma$ can bridge this gap without destroying all useful temporal information.

\begin{table}[h]
\centering
\caption{TCN on Cheetah Run: $\sigma$=0.5 reduces severity 4.7$\times$ but fails to prevent CIVO.}
\label{tab:tcn-cheetah}
\begin{tabular}{lcccc}
\toprule
\textbf{$\sigma$} & \textbf{NVE Max} & \textbf{Eval Peak} & \textbf{Prevention} & \textbf{Note} \\
\midrule
0 (untreated) & 70.6 & 248 & No & 3/3 collapse \\
0.2 & 1.9 & 300 & No ($\sigma$ insufficient) & NVE borderline \\
0.5 & 15.0 & 220 & No (still insufficient) & 4.7$\times$ reduction \\
\bottomrule
\end{tabular}
\end{table}

\paragraph{PickCube failure.}
TCN $\sigma$=0.2 on PickCube (a manipulation task) also fails: NVE peaks at 2.8$\times$ with eval reaching only 35.
PickCube's dynamics---discontinuous contact, sparse reward---differ fundamentally from locomotion, and $\sigma$=0.2 is not calibrated for this task domain.

\subsection{Limitations of TCN}

The TCN results on Walker Walk represent the first complete CIVO prevention in this project, but three limitations constrain the claim:

\begin{enumerate}[leftmargin=*,itemsep=2pt]
\item \textbf{Environment specificity}: Prevention is demonstrated only on Walker Walk. Cheetah Run resists $\sigma$=0.5, and the optimal $\sigma$ appears to be environment-specific with no universal value.
\item \textbf{Narrow effective band}: On Walker Walk, $\sigma$=0.1 is insufficient and $\sigma$=0.5 is destructive. The effective window $\sigma \in [0.15, 0.35]$ is narrow.
\item \textbf{Seed-dependent performance}: Even within the effective $\sigma$, performance varies 50--90\% of baseline across seeds.
\end{enumerate}

We therefore scope C8 as an \emph{exploratory} result: TCN validates the forward-pass intervention principle but requires adaptive calibration mechanisms (beyond this paper's scope) for practical deployment.

%==============================================================================
% 6. DISCUSSION
%==============================================================================
\section{Discussion}

\paragraph{MPC vs.\ Imagination.}
TAWM~\cite{nhu2025tawm} demonstrates that $\dt$-conditioning succeeds in TD-MPC2, an MPC-based architecture.
The critical difference is that MPC re-plans from \emph{real environment states} at each step, while imagination-based actor-critic trains entirely within model rollouts.
In imagination, value estimation errors compound through bootstrapping: the critic learns from its own (overestimated) predictions, creating the positive feedback loop central to CIVO.
MPC's re-grounding in reality breaks this loop.
This architectural distinction---not the temporal conditioning mechanism per se---explains why the same idea succeeds in one paradigm and fails in another.

\paragraph{Intervention Difficulty Gradient.}
Our results establish a hierarchy of intervention difficulty:
\begin{itemize}[leftmargin=*,itemsep=1pt]
\item \textbf{Backward-pass (SGS)}: Categorically fails (1/17). Gradient-level interventions cannot address forward-pass representational damage; compensatory amplification through remaining channels defeats targeted blocking.
\item \textbf{Forward-pass uniform noise (TCN)}: Succeeds on Walker Walk ($\sigma$=0.2) but fails on Cheetah Run. Uniform noise cannot adapt to environment-specific CIVO severity.
\item \textbf{Adaptive forward-pass}: The natural next step---calibrating noise or other forward-pass modifications to local CIVO severity---is beyond this paper's scope but identified as the critical direction.
\end{itemize}

\paragraph{Stochastic Bifurcation.}
CIVO collapse is stochastic, not deterministic.
Among trained \texttt{dt\_embedding} configurations, $\sim$50\% of seeds collapse (across environments).
The NVE\,=\,1.5$\times$ threshold marks a ``point of no return'': when overestimation is sustained above this level for consecutive checkpoints, the feedback loop locks in.
Early training stochasticity (weight initialization, mini-batch sampling) determines which trajectory is taken, consistent with findings on training instability in deep RL~\cite{henderson2018deep}.

\paragraph{Value Self-Correction.}
Seed~2 of 008d demonstrates that transient NVE\,$>$\,1.5$\times$ can self-correct if the overestimation does not persist.
This seed reaches NVE\,=\,2.05$\times$ at 25K but declines to 0.38$\times$ by 105K, achieving eval\,=\,268 (modest but not collapsed).
The self-correction mechanism is not fully understood but appears related to value loss gradients eventually dominating the overestimation drive---provided the positive feedback loop has not yet locked.

\paragraph{Implications for Temporal Abstraction.}
Our diagnostic framework suggests that future temporal abstraction methods in imagination-based RL should:
(1)~avoid training temporal conditioning on multi-step transitions unless the world model's fidelity at those scales is validated;
(2)~use hierarchical architectures (THICK, Hieros) that maintain separate dynamics per temporal scale;
(3)~incorporate forward-pass calibration mechanisms to limit conditioning precision during imagination.

%==============================================================================
% 7. CONCLUSION
%==============================================================================
\section{Conclusion}

We present a systematic diagnostic analysis of why temporal abstraction fails in DreamerV3's imagination-based actor-critic framework.
Through 31+ configurations across four environments, we identify conditioning-induced value overestimation (CIVO) as the root failure mode, with the trained \texttt{dt\_embedding} as its necessary and sufficient trigger.
The NVE\,=\,1.5$\times$ diagnostic threshold achieves AUROC\,=\,0.969 for detecting pathological overestimation.
Gradient-level interventions (SGS) fail comprehensively (1/17 seeds), establishing that backward-pass approaches cannot address the forward-pass vulnerability.
As an exploratory result, TCN ($\sigma$=0.2) achieves the first complete CIVO prevention on Walker Walk, but the effective noise band is narrow and environment-specific, pointing toward adaptive forward-pass interventions as the next research direction.

Our findings carry a cautionary message for the temporal abstraction research agenda in world-model RL: the mechanism that enables precise temporal conditioning is the same mechanism that creates the exploitation vulnerability.
Resolving this tension---maintaining temporal awareness without enabling CIVO---is the core challenge for future work.

%==============================================================================
% REFERENCES
%==============================================================================
\bibliographystyle{plainnat}

\begin{thebibliography}{20}

\bibitem[Hafner et~al.(2023)]{hafner2023dreamerv3}
D.~Hafner, J.~Pasukonis, J.~Ba, and T.~Lillicrap.
\newblock Mastering diverse domains through world models.
\newblock \emph{arXiv preprint arXiv:2301.04104}, 2023.

\bibitem[Hafner et~al.(2021)]{hafner2021dreamerv2}
D.~Hafner, T.~Lillicrap, M.~Norouzi, and J.~Ba.
\newblock Mastering {A}tari with discrete world models.
\newblock In \emph{ICLR}, 2021.

\bibitem[Hafner et~al.(2020)]{hafner2020dream}
D.~Hafner, T.~Lillicrap, J.~Ba, and M.~Norouzi.
\newblock Dream to control: Learning behaviors by latent imagination.
\newblock In \emph{ICLR}, 2020.

\bibitem[Sutton et~al.(1999)]{sutton1999options}
R.~S.~Sutton, D.~Precup, and S.~Singh.
\newblock Between {MDP}s and semi-{MDP}s: A framework for temporal abstraction in reinforcement learning.
\newblock \emph{Artificial Intelligence}, 112(1--2):181--211, 1999.

\bibitem[Precup(2000)]{precup2000temporal}
D.~Precup.
\newblock Temporal abstraction in reinforcement learning.
\newblock PhD thesis, University of Massachusetts Amherst, 2000.

\bibitem[Lakshminarayanan et~al.(2017)]{lakshminarayanan2017dynamic}
A.~S.~Lakshminarayanan, S.~Sharma, and B.~Ravindran.
\newblock Dynamic action repetition for deep reinforcement learning.
\newblock In \emph{AAAI}, 2017.

\bibitem[Sharma et~al.(2017)]{sharma2017figar}
S.~Sharma, A.~S.~Lakshminarayanan, and B.~Ravindran.
\newblock Fine grained action repetition for deep reinforcement learning.
\newblock \emph{arXiv preprint arXiv:1702.06054}, 2017.

\bibitem[Nhu et~al.(2025)]{nhu2025tawm}
V.~Nhu, et~al.
\newblock {TAWM}: Temporal abstraction in world models.
\newblock In \emph{ICML}, 2025.

\bibitem[Kim et~al.(2023)]{kim2023thick}
S.~Kim, S.~Woo, and D.~Kim.
\newblock {THICK}: Temporal hierarchy for context in world models.
\newblock In \emph{NeurIPS}, 2023.

\bibitem[{Hieros}(2024)]{hieros2024}
Hieros: Hierarchical imagination on structured state space sequence world models.
\newblock In \emph{ICML}, 2024.

\bibitem[Janner et~al.(2019)]{janner2019mbpo}
M.~Janner, J.~Fu, M.~Zhang, and S.~Levine.
\newblock When to trust your model: Model-based policy optimization.
\newblock In \emph{NeurIPS}, 2019.

\bibitem[Lu et~al.(2018)]{lu2018delusional}
T.~Lu, D.~Schuurmans, and C.~Boutilier.
\newblock Non-delusional {Q}-learning and value-iteration.
\newblock In \emph{NeurIPS}, 2018.

\bibitem[Venkatraman et~al.(2015)]{venkatraman2015improving}
A.~Venkatraman, M.~Hebert, and J.~A.~Bagnell.
\newblock Improving multi-step prediction of learned time series models.
\newblock In \emph{AAAI}, 2015.

\bibitem[Henderson et~al.(2018)]{henderson2018deep}
P.~Henderson, R.~Islam, P.~Bachman, J.~Pineau, D.~Precup, and D.~Meger.
\newblock Deep reinforcement learning that matters.
\newblock In \emph{AAAI}, 2018.

\bibitem[Agarwal et~al.(2021)]{agarwal2021deep}
R.~Agarwal, M.~Schwarzer, P.~S.~Castro, A.~C.~Courville, and M.~G.~Bellemare.
\newblock Deep reinforcement learning at the edge of the statistical precipice.
\newblock In \emph{NeurIPS}, 2021.

\end{thebibliography}

%==============================================================================
% APPENDIX
%==============================================================================
\appendix
\section{Complete Experimental Results}

\subsection{Walker Walk 008d Evaluation Trajectories}

\begin{table}[h]
\centering
\caption{Walker Walk 008d (trained \texttt{dt\_emb}, $\dt$=2, subsample ON, H=8): full evaluation trajectories.}
\label{tab:appendix-008d-traj}
\begin{tabular}{lcc}
\toprule
\textbf{Step} & \textbf{Seed 1 Eval} & \textbf{Seed 2 Eval} \\
\midrule
5K & 31 & 39 \\
15K & 155 & 144 \\
25K & 170 & 141 \\
35K & 150 & 164 \\
45K & 143 & 141 \\
55K & 85 & 172 \\
65K & 245 & 209 \\
75K & 180 & 197 \\
85K & 200 & 236 \\
95K & 135 & 279 \\
105K & 130 & 268 \\
\bottomrule
\end{tabular}
\end{table}

\subsection{Walker Walk 008f Evaluation Trajectories}

\begin{table}[h]
\centering
\caption{Walker Walk 008f (\texttt{dt\_embed\_dim=0}, subsample ON): full evaluation trajectories (seeds 1--3).}
\label{tab:appendix-008f-traj}
\begin{tabular}{lccc}
\toprule
\textbf{Step} & \textbf{Seed 1 Eval} & \textbf{Seed 2 Eval} & \textbf{Seed 3 Eval} \\
\midrule
5K & 30 & 30 & 31 \\
15K & 165 & 157 & 148 \\
25K & 195 & 169 & 180 \\
35K & 210 & 179 & 220 \\
45K & 190 & 211 & 255 \\
55K & 175 & 239 & 300 \\
65K & 370 & 244 & 280 \\
75K & 320 & 359 & 310 \\
85K & 250 & 227 & 350 \\
95K & 320 & 191 & 290 \\
105K & 500 & 148 & 312 \\
\bottomrule
\end{tabular}
\end{table}

\subsection{Walker Walk 008e\_v3 Evaluation Trajectories}

\begin{table}[h]
\centering
\caption{Walker Walk 008e\_v3 ($\dt$=1 imagination, trained \texttt{dt\_emb}): evaluation trajectories (4 seeds).}
\label{tab:appendix-008ev3-traj}
\begin{tabular}{lcccc}
\toprule
\textbf{Step} & \textbf{Seed 1} & \textbf{Seed 2} & \textbf{Seed 3} & \textbf{Seed 4} \\
\midrule
5K & 31 & 39 & 35 & 38 \\
15K & 131 & 181 & 145 & 160 \\
25K & 156 & 212 & 195 & 200 \\
35K & 276 & 320 & 310 & 330 \\
45K & 382 & 357 & 380 & 420 \\
55K & 487 & 418 & 440 & 510 \\
65K & 581 & 435 & 460 & 600 \\
75K & 574 & 577 & 510 & 640 \\
85K & 657 & 715 & 520 & 720 \\
95K & 852 & 452 & 320 & 830 \\
105K & 870 & 137 & 93 & 899 \\
\bottomrule
\end{tabular}
\end{table}

\subsection{Cheetah Run Trained \texttt{dt\_emb} Trajectories}

\begin{table}[h]
\centering
\caption{Cheetah Run trained \texttt{dt\_emb} (subsample ON): evaluation trajectories (3 seeds). All collapse.}
\label{tab:appendix-cheetah-traj}
\begin{tabular}{lccc}
\toprule
\textbf{Step} & \textbf{Seed 0} & \textbf{Seed 1} & \textbf{Seed 2} \\
\midrule
5K & 8 & 10 & 9 \\
15K & 80 & 95 & 70 \\
25K & 180 & 220 & 150 \\
35K & 248 & 310 & 230 \\
45K & 200 & 280 & 150 \\
55K & 120 & 170 & 60 \\
65K & 45 & 80 & 25 \\
75K & 20 & 30 & 10 \\
85K & 5 & 15 & 3 \\
95K & 1 & 52 & 1 \\
105K & 1 & 52 & 1 \\
\bottomrule
\end{tabular}
\end{table}

\subsection{Frozen \texttt{dt\_embedding} Results}

\begin{table}[h]
\centering
\caption{Walker Walk frozen \texttt{dt\_emb} (015b): prevents CIVO but suffers gradient starvation.}
\label{tab:appendix-frozen}
\begin{tabular}{lcccccc}
\toprule
\textbf{Step} & 5K & 25K & 45K & 65K & 85K & 105K \\
\midrule
Eval & 31 & 172 & 122 & 160 & 184 & 204 \\
NVE & --- & --- & --- & --- & 0.47 & 0.42 \\
\bottomrule
\end{tabular}
\end{table}

\subsection{TCN Walker s0 $\sigma$=0.1 Projected Trajectory}

\begin{table}[h]
\centering
\caption{TCN Walker Walk $\sigma$=0.1 (s0): NVE briefly breaches 1.5$\times$ at 65K---insufficient noise.}
\label{tab:appendix-tcn-01}
\begin{tabular}{lcc}
\toprule
\textbf{Step} & \textbf{Eval} & \textbf{NVE} \\
\midrule
25K & 200 & --- \\
35K & 280 & --- \\
55K & 300 & --- \\
65K & 250 & 1.65 \\
85K & 280 & --- \\
105K & 380 & --- \\
135K & 500 & --- \\
175K & 580 & --- \\
215K & 550 & 0.25 \\
\bottomrule
\end{tabular}
\end{table}

\subsection{TCN Walker s1 $\sigma$=0.2 Projected Trajectory}

\begin{table}[h]
\centering
\caption{TCN Walker Walk $\sigma$=0.2 seed 1: CIVO prevented (NVE\,$<$\,1.0$\times$) but eval peak 460 (50\% baseline).}
\label{tab:appendix-tcn-02-s1}
\begin{tabular}{lcc}
\toprule
\textbf{Step} & \textbf{Eval} & \textbf{NVE} \\
\midrule
85K & 152 & 0.40 \\
105K & 210 & 0.65 \\
125K & 310 & 0.42 \\
145K & 380 & 0.32 \\
175K & 460 & 0.22 \\
195K & 420 & 0.20 \\
215K & 445 & 0.18 \\
\bottomrule
\end{tabular}
\end{table}

\subsection{TCN Cheetah $\sigma$=0.5 Projected Trajectory}

\begin{table}[h]
\centering
\caption{TCN Cheetah Run $\sigma$=0.5 (s0): severity reduced 4.7$\times$ vs.\ untreated but still pathological.}
\label{tab:appendix-tcn-cheetah}
\begin{tabular}{lcc}
\toprule
\textbf{Step} & \textbf{Eval} & \textbf{NVE} \\
\midrule
15K & 78 & --- \\
35K & 220 & --- \\
55K & 130 & 4.5 \\
65K & 80 & 8.2 \\
85K & 30 & --- \\
105K & 10 & 15.0 \\
\bottomrule
\end{tabular}
\end{table}

\subsection{TCN PickCube $\sigma$=0.2}

TCN $\sigma$=0.2 on the PickCube manipulation task yields eval peak 35, NVE peak 2.8$\times$, and final eval 22.
The pattern is transient CIVO followed by value collapse, similar to untreated PickCube behavior.
PickCube's discontinuous contact dynamics and sparse reward make it fundamentally different from locomotion tasks; $\sigma$=0.2 is not calibrated for this domain.

\end{document}
